% Directly inputtable article fragment.
% Required preamble packages: amsmath, amssymb.
% Author: ChatGPT, 2 September 2026.

\section{Refined tensor Haar volume and the pointed same-pilot gate}
\label{sec:refined-tensor-haar-same-pilot}

\noindent\textbf{Author: ChatGPT.}\quad
This section isolates an unconditional local theorem that is useful at the
present Project-LANA/IUT interface.  It does not assert IUT III,
Corollary~3.12, or the $abc$ conjecture.  In particular, the finite-etale
factorization, Haar-positivity, Ind3-openness, and pointed-horizontal
identification hypotheses below are logically distinct.

\noindent\textbf{Primary-source boundary.}
The terminology follows Mochizuki's May 2020 author version of
\emph{Inter-universal Teichmuller Theory III}.  The exact locators are
Proposition~3.9(i), Remark~3.9.5(i), Theorem~3.11(i)--(ii), and
Corollary~3.12, especially proof Steps (xi-b)--(xi-f).  The identification
still missing after the three indeterminacies is also recorded in Sections
6--7 of the July 2026 LANA formalization report.  The frozen Project LANA
source is commit
\begin{center}
\texttt{c65b28c9f9631635e742294c3a5df15759e7c74c}.
\end{center}
No source claim is imported as a premise unless it is displayed below.

\subsection{Primitive factors and the prime Jacobian}

Fix a rational prime $p$.  For a place tuple $c$, let
\[
 A_c=\bigotimes_{j\in S,\,\mathbb Q_p}K_{c_j}.
\]
Since the local extensions are finite and separable, $A_c$ is a finite-etale
$\mathbb Q_p$-algebra.  Thus
\begin{equation}
 A_c\simeq\prod_{d\in D_c}L_{c,d},
 \qquad
 N_c:=\dim_{\mathbb Q_p}A_c
      =\sum_{d\in D_c}[L_{c,d}:\mathbb Q_p].
 \label{eq:primitive-decomposition}
\end{equation}
For a factor $L=L_{c,d}$, write $e=e(L/\mathbb Q_p)$ and
$f=f(L/\mathbb Q_p)$.

\paragraph{Primitive-factor theorem.}
For every uniformizer $\pi\in\mathcal O_L$, there is
$u\in\mathcal O_L^\times$ such that
\begin{equation}
 p=u\pi^e,\qquad |k_L|=p^f,
 \qquad [L:\mathbb Q_p]=ef.
 \label{eq:ef-identities}
\end{equation}
Indeed, unique factorization in the discrete valuation ring gives
$p=u\pi^m$.  Applying the integer-valued valuation normalized by
$v_L(\pi)=1$ yields $m=v_L(p)=e$.  The residue field is an
$f$-dimensional vector space over $\mathbb F_p$, and the fundamental equality
for finite extensions of complete discretely valued fields gives
$[L:\mathbb Q_p]=ef$.  Consequently,
\begin{equation}
 N_c=\sum_{d\in D_c}e_{c,d}f_{c,d}>0.
 \label{eq:dimension-conservation}
\end{equation}
The last identity follows by taking dimensions in
\eqref{eq:primitive-decomposition}; it is not an assumption about an
unsplit tensor product.

Normalize additive Haar measure $\mu_{c,d}$ by
$\mu_{c,d}(\mathcal O_{c,d})=1$.  For nonempty compact-open
$U_{c,d}\subseteq L_{c,d}$, set
\[
 U_c=\prod_dU_{c,d},\qquad
 \mu_c=\prod_d\mu_{c,d},\qquad
 \lambda_c(U_c)=\frac{1}{N_c}\log\mu_c(U_c).
\]
All measures in this display are finite and strictly positive.  The additive
Haar change-of-variables formula and \eqref{eq:ef-identities} give
\[
 \mu_{c,d}(p^{-1}U_{c,d})
   =p^{e_{c,d}f_{c,d}}\mu_{c,d}(U_{c,d}).
\]
Taking products and using \eqref{eq:dimension-conservation} proves
\begin{equation}
 \lambda_c(p^{-1}U_c)=\lambda_c(U_c)+\log p.
 \label{eq:normalized-prime-shift}
\end{equation}
Moreover, with $n_{c,d}=[L_{c,d}:\mathbb Q_p]$ and
$\rho_{c,d}=n_{c,d}/N_c$, one has
\begin{equation}
 \sum_d\rho_{c,d}=1,
 \qquad
 \lambda_c(U_c)=
 \sum_d\rho_{c,d}
   \frac{\log\mu_{c,d}(U_{c,d})}{n_{c,d}}.
 \label{eq:relative-degree-average}
\end{equation}
Thus the normalization is by total local degree, rather than by the number
of primitive factors.

If a tuple has weight $W_c$, its primitive factors must receive
\begin{equation}
 \widetilde W_{c,d}=W_c\rho_{c,d}.
 \label{eq:refined-weights}
\end{equation}
Then $\sum_d\widetilde W_{c,d}=W_c$, so normalized tuple weights remain
normalized after primitive refinement.

\subsection{Finite-positive orbit regions}

Let $X$ be a finite product of locally compact primitive fields with product
Haar measure $\mu$.  Let $U\subseteq X$ be nonempty and open, and let
$\{\phi_a:X\simeq X\}_{a\in A}$ be homeomorphisms, one of which is the
identity.  Suppose that every $\phi_a(U)$ lies in a common set $E$ with
compact closure.  Then
\begin{equation}
 \Theta(U):=\bigcup_{a\in A}\phi_a(U)
 \label{eq:theta-orbit}
\end{equation}
is nonempty and open.  It therefore has positive Haar measure.  Its closure
is a closed subset of the compact set $\overline E$, so it is compact; Haar
measure is finite on that closure.  Hence
\begin{equation}
 0<\mu(\Theta(U))<\infty.
 \label{eq:theta-finite-positive}
\end{equation}
When $A$ is finite and $U$ is compact-open, $\Theta(U)$ is itself
compact-open.  For an infinite family the stated hypotheses yield an open
region with compact closure; they do not imply that the orbit union is
closed.

The same argument applies to an inductively generated Ind1--Ind3 system only
under the following separate premises: ordinary branches are open and one is
nonempty; Ind1 and Ind2 preserve openness; every relational Ind3 target of an
open reachable region is open; and all constructors preserve the common
compact-envelope containment.  Structural induction then proves openness
and envelope containment for every reachable region, after which
\eqref{eq:theta-finite-positive} follows.  The finite-etale theorem above
does not supply any of these source-level Ind hypotheses.

\subsection{The pointed same-pilot implication}

Let $P$ be the actual native $q$-pilot region, $\Theta$ the actual union of
reachable theta outputs, and $H$ a finite-positive hull.  If the
object-level construction proves the pointed inclusions
\begin{equation}
 P\subseteq\Theta\subseteq H,
 \label{eq:pointed-inclusions}
\end{equation}
then Haar monotonicity and monotonicity of the real logarithm give
\begin{equation}
 \log\mu(P)\leq\log\mu(\Theta)\leq\log\mu(H).
 \label{eq:same-pilot-sandwich}
\end{equation}
Thus any independent estimate $\log\mu(H)\leq T$ implies
$\log\mu(P)\leq T$.  No volume-preservation assertion for the Ind maps is
needed for this conditional implication.  However, naming some generated
branch ``native'' does not establish \eqref{eq:pointed-inclusions}.  The open
source-level obligation is to transport the original pointed pilot through
the theta link, IPL/SHE/APT, log-Kummer correction, determinant
normalization, and Ind3.

\subsection{Exact counterexamples and route status}

Each counterexample below satisfies all premises of the deliberately weaker
claim that it refutes.

\begin{enumerate}
\item A one-factor quadratic local field has factor count $1$ and total
degree $2$.  Its raw prime-preimage shift is $2\log p$, so division by factor
count does not give \eqref{eq:normalized-prime-shift}.  This refutes only
factor-count normalization.
\item For an unramified quadratic Galois extension $K/\mathbb Q_p$,
$K\otimes_{\mathbb Q_p}K\simeq K\times K$.  Copying a parent weight $1$ to
both factors gives total weight $2$; the relative-degree weights
$(1/2,1/2)$ satisfy \eqref{eq:refined-weights}.  This refutes only copied
parent weights.
\item On $\mathbb R$ with Lebesgue Haar measure, an Ind3 relation may send
$(-1,1)$ to $\{0\}$ while both lie in the compact envelope $[-1,1]$.
The target has measure zero.  This refutes the claim that envelope
preservation alone yields finite-positive Ind3 outputs; it does not refute a
concrete Ind3 theorem that proves openness or positivity.
\item In the discrete three-point group with counting Haar measure, take
$P=\{0,1\}$ and $\Theta=H=\{0\}$.  All regions are compact-open and
finite-positive and $\Theta\subseteq H$, but
$\log\mu(P)=\log2>0=\log\mu(H)$.  This refutes only an unpointed inference.
\end{enumerate}

The finite-etale factor identities and total-degree Haar normalization are
proved.  The relative-degree weights, bounded-open orbit lemma, and
conditional pointed sandwich are also proved.  The identification of the
concrete IUT Ind1 and Ind2 maps,
openness or positive measure of the required Ind3 targets, and the pointed
IPL, SHE, and APT horizontal transport remain open gaps.  No counterexample
satisfying all premises of those source statements is presently supplied;
therefore those routes remain active.  IUT III, Corollary~3.12, IUT, and the
standard $abc$ conjecture remain open in this work.

\subsection{Formalization boundary}

The companion Lean development derives the identities in
\eqref{eq:ef-identities} and \eqref{eq:dimension-conservation} on the actual
primitive quotients of a \texttt{TupleFiniteEtalePacket}.  It also formalizes
the product-Haar formulas, the refined weights, the orbit and Ind interface
lemmas and the pointed sandwich.  It checks coefficient-level versions of the
factor-count and copied-weight obstructions and set-level versions of the
envelope and unpointed obstructions.  The quadratic local-field and tensor
decomposition witnesses for the first two remain paper proofs.  Its axiom
audit reports only \texttt{propext}, \texttt{Classical.choice}, and
\texttt{Quot.sound}; there are no project axioms or admitted proofs.

The following source-specific work remains on paper: identifying each IUT
capsule's completed tensor algebra with the finite-etale packet; comparing
the tensor product of integral orders with the product of maximal factor
orders, including finite-index/different corrections; instantiating the
abstract Ind1/Ind2 homeomorphisms and common envelope; proving the required
Ind3 openness or positivity; and constructing the pointed horizontal map.
Finite computation audits the rational normalization formulas only and gives
no evidence for these source-specific obligations or for $abc$.
